% =============================================================================
\chapter{RL avec Contraintes et S\'ecurit\'e}
\label{ch:contraintes}
% =============================================================================

Un robot qui apprend \`a marcher peut tomber. Un v\'ehicule autonome qui apprend \`a conduire peut provoquer un accident. Un algorithme de trading qui apprend \`a investir peut faire faillite. Dans ces applications, maximiser la r\'ecompense ne suffit pas~: il faut \emph{garantir la s\'ecurit\'e}, y compris pendant la phase d'apprentissage. Le RL contraint (Constrained MDP, Altman 1999) formalise ce probl\`eme en ajoutant des contraintes sur les co\^uts cumulatifs~; le RL s\^ur (Safe RL) vise des garanties plus fortes, \`a chaque pas de temps. Les m\'ethodes lagrangiennes, les m\'ethodes \`a barri\`eres, et les approches fond\'ees sur les fonctions de Lyapunov (Chow et al., 2018) forment l'arsenal de ce domaine en pleine expansion.

\begin{intuition}
Dans de nombreuses applications r\'eelles (robotique, v\'ehicules autonomes,
sant\'e), maximiser la r\'ecompense ne suffit pas~: il faut \'egalement
respecter des contraintes de s\'ecurit\'e. Le RL contraint (\emph{Constrained RL})
formalise ce probl\`eme en ajoutant des contraintes sur les co\^uts cumulatifs.
Le RL s\^ur (\emph{Safe RL}) vise \`a garantir la s\'ecurit\'e \`a chaque instant,
y compris pendant l'apprentissage.
\end{intuition}

% -----------------------------------------------------------------------------
\section{MDP contraint (CMDP)}
% -----------------------------------------------------------------------------

\begin{definition}[CMDP]
Un \textbf{MDP contraint} est un MDP augment\'e de $K$ fonctions de co\^ut
$c_k : \mathcal{S} \times \mathcal{A} \to \R$ et de seuils $d_k$~:
\[
    \max_\pi J_R(\pi) = \E_\pi\left[\sum_{t=0}^\infty \gamma^t R_{t+1}\right]
    \quad \text{s.c.} \quad
    J_{c_k}(\pi) = \E_\pi\left[\sum_{t=0}^\infty \gamma^t c_k(S_t, A_t)\right] \leq d_k
    \quad \forall k.
\]
\end{definition}

\begin{theoreme}[Existence d'une politique optimale pour CMDP]
Pour un CMDP fini avec $K$ contraintes, il existe une politique optimale
qui est \textbf{m\'elange} d'au plus $K+1$ politiques d\'eterministes
stationnaires. Si l'ensemble r\'ealisable a un int\'erieur non vide (condition
de Slater), le probl\`eme peut \^etre r\'esolu par la m\'ethode du Lagrangien.
\end{theoreme}

% -----------------------------------------------------------------------------
\section{Approche lagrangienne}
% -----------------------------------------------------------------------------

\begin{definition}[Lagrangien du CMDP]
Le \textbf{Lagrangien} du CMDP est~:
\[
    \mathcal{L}(\pi, \boldsymbol{\lambda}) = J_R(\pi)
    - \sum_{k=1}^K \lambda_k \left(J_{c_k}(\pi) - d_k\right)
\]
o\`u $\boldsymbol{\lambda} = (\lambda_1, \ldots, \lambda_K) \geq 0$ sont les
multiplicateurs de Lagrange.
\end{definition}

\begin{theoreme}[Dualit\'e forte]
Sous la condition de Slater (il existe $\pi_0$ tel que $J_{c_k}(\pi_0) < d_k$
pour tout $k$), la dualit\'e forte s'applique~:
\[
    \max_\pi \min_{\boldsymbol{\lambda} \geq 0} \mathcal{L}(\pi, \boldsymbol{\lambda})
    = \min_{\boldsymbol{\lambda} \geq 0} \max_\pi \mathcal{L}(\pi, \boldsymbol{\lambda}).
\]
Le probl\`eme dual est convexe en $\boldsymbol{\lambda}$.
\end{theoreme}

\begin{algorithme}[title=\textbf{PPO-Lagrangien}]
\begin{enumerate}
    \item Initialiser $\theta$ (politique), $\phi$ (critique de r\'ecompense),
    $\psi$ (critique de co\^ut), $\lambda \geq 0$
    \item Pour chaque it\'eration~:
    \begin{enumerate}
        \item Collecter des trajectoires avec $\pi_\theta$
        \item Estimer les avantages de r\'ecompense $\hat{A}_t^R$ et de co\^ut $\hat{A}_t^C$ via GAE
        \item Mettre \`a jour $\theta$ par PPO-Clip sur l'objectif lagrangien~:
        \[
            L(\theta) = L_R^{\text{CLIP}}(\theta)
            - \lambda \, L_C^{\text{CLIP}}(\theta)
        \]
        \item Mettre \`a jour $\lambda$~:
        $\lambda \leftarrow \max(0, \lambda + \eta_\lambda(J_C(\pi_\theta) - d))$
        \item Mettre \`a jour les critiques $\phi$ et $\psi$
    \end{enumerate}
\end{enumerate}
\end{algorithme}

% -----------------------------------------------------------------------------
\section{CPO --- Constrained Policy Optimization}
% -----------------------------------------------------------------------------

\begin{definition}[CPO]
\textbf{CPO} (Achiam et al., 2017) \'etend TRPO au cadre contraint en r\'esolvant~:
\[
    \max_\theta \; \E_t[r_t(\theta) \hat{A}_t^R] \quad \text{s.c.} \quad
    \begin{cases}
        J_{c_k}(\pi_{\theta_{\text{old}}}) + \E_t[r_t(\theta) \hat{A}_t^{c_k}] \leq d_k & \forall k \\
        \text{KL}[\pi_{\theta_{\text{old}}} \| \pi_\theta] \leq \delta
    \end{cases}
\]
Ce probl\`eme quadratique sous contraintes lin\'eaires et quadratiques se r\'esout
analytiquement.
\end{definition}

\begin{theoreme}[Garantie de CPO]
Sous certaines conditions (lin\'earisation valide), CPO garantit la satisfaction
des contraintes de co\^ut \`a chaque mise \`a jour de politique, pas seulement
\`a convergence.
\end{theoreme}

% -----------------------------------------------------------------------------
\section{Approches par barri\`ere et p\'enalit\'e}
% -----------------------------------------------------------------------------

\begin{definition}[M\'ethode de barri\`ere logarithmique]
On remplace la contrainte par une p\'enalit\'e barri\`ere~:
\[
    \max_\pi J_R(\pi) - \frac{1}{t} \sum_k \log(d_k - J_{c_k}(\pi))
\]
o\`u $t > 0$ contr\^ole la pr\'ecision de l'approximation. Quand $t \to \infty$,
la solution approche celle du probl\`eme contraint.
\end{definition}

\begin{definition}[P\'enalit\'e quadratique augment\'ee]
La m\'ethode du \textbf{Lagrangien augment\'e} combine multiplicateurs et p\'enalit\'e~:
\[
    \mathcal{L}_\rho(\pi, \lambda) = J_R(\pi) - \lambda(J_C(\pi) - d)
    - \frac{\rho}{2}(J_C(\pi) - d)_+^2
\]
o\`u $(x)_+ = \max(0, x)$ et $\rho > 0$ est le param\`etre de p\'enalit\'e.
\end{definition}

% -----------------------------------------------------------------------------
\section{Safe RL --- S\'ecurit\'e pendant l'apprentissage}
% -----------------------------------------------------------------------------

\begin{definition}[S\'ecurit\'e au sens de Lyapunov]
Un \'etat $s$ est \textbf{s\^ur} s'il existe une politique $\pi$ qui maintient
le syst\`eme dans un ensemble s\^ur $\mathcal{S}_{\text{safe}} \subseteq \mathcal{S}$
pour tout temps futur. La \textbf{fonction de Lyapunov} $L(s) \geq 0$ satisfait~:
\[
    \E[L(S_{t+1}) \mid S_t = s, A_t = a] - L(s) \leq -\epsilon
    \quad \forall s \notin \mathcal{S}_{\text{goal}}.
\]
\end{definition}

\begin{definition}[Bouclier de s\'ecurit\'e (\emph{Safety Shield})]
Un \textbf{bouclier de s\'ecurit\'e} est un module qui filtre les actions
propos\'ees par l'agent RL~:
\[
    a_{\text{safe}} = \begin{cases}
        a & \text{si } a \in \mathcal{A}_{\text{safe}}(s) \\
        \argmin_{a' \in \mathcal{A}_{\text{safe}}(s)} \|a' - a\| & \text{sinon}
    \end{cases}
\]
o\`u $\mathcal{A}_{\text{safe}}(s) \subseteq \mathcal{A}$ est l'ensemble des
actions s\^ures dans l'\'etat $s$.
\end{definition}

% -----------------------------------------------------------------------------
\section{Fonctions barri\`ere de contr\^ole (CBF)}
% -----------------------------------------------------------------------------

\begin{definition}[Control Barrier Function]
Une \textbf{CBF} $h : \mathcal{S} \to \R$ d\'efinit l'ensemble s\^ur
$\mathcal{C} = \{s : h(s) \geq 0\}$. Pour un syst\`eme $s_{t+1} = f(s_t, a_t)$,
la condition de s\'ecurit\'e est~:
\[
    h(f(s, a)) - h(s) \geq -\kappa \, h(s), \quad \kappa \in (0, 1).
\]
Si cette condition est satisfaite, $h(s_t) \geq 0$ pour tout $t$ si $h(s_0) \geq 0$.
\end{definition}

\begin{algorithme}[title=\textbf{RL avec filtre CBF}]
\begin{enumerate}
    \item L'agent RL propose une action $a_{\text{rl}}$
    \item R\'esoudre le QP (programme quadratique)~:
    \[
        a^* = \argmin_a \|a - a_{\text{rl}}\|^2 \quad \text{s.c.} \quad
        h(f(s, a)) - h(s) \geq -\kappa \, h(s)
    \]
    \item Ex\'ecuter $a^*$ dans l'environnement
\end{enumerate}
\end{algorithme}

% -----------------------------------------------------------------------------
\section{Impl\'ementation}
% -----------------------------------------------------------------------------

\begin{codeblock}[title=\textbf{PPO-Lagrangien simplifi\'e}]
\begin{minted}{python}
import torch
import torch.nn as nn
import numpy as np

class LagrangianPPO:
    def __init__(self, actor_critic, cost_critic, cost_limit=25.0,
                 lambda_lr=0.01, clip_eps=0.2, gamma=0.99):
        self.ac = actor_critic
        self.cost_critic = cost_critic
        self.cost_limit = cost_limit
        self.log_lambda = torch.tensor(0.0, requires_grad=True)
        self.lambda_opt = torch.optim.Adam([self.log_lambda], lr=lambda_lr)
        self.clip_eps = clip_eps
        self.gamma = gamma

    @property
    def lam(self):
        return self.log_lambda.exp().item()

    def compute_policy_loss(self, states, actions, old_log_probs,
                            reward_advantages, cost_advantages):
        dist, _ = self.ac(states)
        new_log_probs = dist.log_prob(actions)
        ratio = (new_log_probs - old_log_probs).exp()

        # Objectif PPO-Clip pour la recompense
        surr1 = ratio * reward_advantages
        surr2 = ratio.clamp(1 - self.clip_eps, 1 + self.clip_eps) * reward_advantages
        reward_loss = -torch.min(surr1, surr2).mean()

        # Objectif PPO pour le cout (on veut le minimiser)
        cost_loss = (ratio * cost_advantages).mean()

        # Objectif lagrangien
        total_loss = reward_loss + self.lam * cost_loss
        return total_loss, dist.entropy().mean()

    def update_lambda(self, mean_episode_cost):
        """Mise a jour du multiplicateur de Lagrange."""
        lambda_loss = -self.log_lambda.exp() * (mean_episode_cost - self.cost_limit)
        self.lambda_opt.zero_grad()
        lambda_loss.backward()
        self.lambda_opt.step()
\end{minted}
\end{codeblock}

\begin{codeblock}[title=\textbf{Filtre de s\'ecurit\'e simple}]
\begin{minted}{python}
import numpy as np
from scipy.optimize import minimize

def safety_filter(action_rl, state, barrier_fn, dynamics_fn,
                  kappa=0.1):
    """
    Filtre l'action pour satisfaire la contrainte CBF.
    barrier_fn: h(s) -> float (positif = sur)
    dynamics_fn: f(s, a) -> s_next
    """
    h_current = barrier_fn(state)

    def constraint(a):
        s_next = dynamics_fn(state, a)
        h_next = barrier_fn(s_next)
        return h_next - (1 - kappa) * h_current

    result = minimize(
        fun=lambda a: np.sum((a - action_rl) ** 2),
        x0=action_rl,
        method='SLSQP',
        constraints={'type': 'ineq', 'fun': constraint},
        bounds=[(-1, 1)] * len(action_rl)
    )
    return result.x if result.success else action_rl
\end{minted}
\end{codeblock}

% -----------------------------------------------------------------------------
\section{Exercices}
% -----------------------------------------------------------------------------

\begin{exercice}[CMDP lagrangien]
Formulez le probl\`eme de navigation s\^ure (un robot doit atteindre un objectif
en \'evitant des zones dangereuses) comme un CMDP. Impl\'ementez PPO-Lagrangien
et v\'erifiez que le co\^ut cumul\'e reste sous le seuil apr\`es convergence.
\end{exercice}

\begin{exercice}[CPO]
Impl\'ementez CPO et comparez-le avec PPO-Lagrangien sur l'environnement
\texttt{SafetyPointGoal-v0} de Safety Gymnasium. Lequel respecte mieux les
contraintes pendant l'entra\^inement~?
\end{exercice}

\begin{exercice}[CBF apprise]
Entra\^inez un r\'eseau de neurones pour apprendre la fonction barri\`ere $h_\psi(s)$
\`a partir de donn\'ees \'etiquet\'ees (\'etats s\^urs vs dangereux). Utilisez cette
CBF comme filtre de s\'ecurit\'e et \'evaluez le taux de violations de contraintes.
\end{exercice}

\begin{exercice}[Th\'eorie]
D\'emontrez que sous la condition de Slater, la dualit\'e forte s'applique
au CMDP et que le probl\`eme dual est convexe. Donnez la forme explicite
du sous-gradient de la fonction duale.
\end{exercice}
